% !TEX program = xelatex
\documentclass[12pt,a4paper,UTF8]{ctexart}

\usepackage[top=2.5cm, bottom=2.5cm, left=3cm, right=2.5cm]{geometry}
\usepackage{fontspec}
\usepackage{xcolor}
\usepackage{hyperref}
\usepackage{amssymb}
\usepackage{booktabs}
\usepackage{longtable}
\usepackage{array}
\usepackage{tabularx}
\usepackage{multirow}
\usepackage{fancyhdr}
\usepackage{titlesec}
\usepackage{caption}
\usepackage{graphicx}
\usepackage{listings}
\usepackage{float}

\hypersetup{colorlinks=true,linkcolor=black,citecolor=black,urlcolor=blue}

\pagestyle{fancy}
\fancyhf{}
\fancyhead[C]{团队工作日报}
\fancyfoot[C]{\thepage}
\renewcommand{\headrulewidth}{0.4pt}

\titleformat{\section}{\large\bfseries}{\thesection}{1em}{}
\titleformat{\subsection}{\normalsize\bfseries}{\thesubsection}{1em}{}
\titleformat{\subsubsection}{\small\bfseries}{\thesubsubsection}{1em}{}

\begin{document}

\begin{titlepage}
\centering
\vspace*{3cm}
{\Huge\bfseries 转转校园二手交易平台\\[0.5cm]}
{\LARGE 团队工作日报}\\[1.5cm]
{\large \textbf{项目周期：} 2026年7月4日 — 2026年7月10日}\\
{\large \textbf{项目名称：} 校园二手物品交易平台}\\[0.3cm]
{\large \textbf{小组成员：} 张可凡（队长）、阮崇轩、刘轩霆、胡欣悦、李硕}\\[5cm]
\vfill
{\large \textbf{编写日期：} 2026年7月10日}
\end{titlepage}

\tableofcontents
\newpage

\section{成员角色分工}

\begin{table}[H]
\centering
\caption{小组成员角色分工}
\begin{tabularx}{\textwidth}{c c X}
\toprule
\textbf{成员} & \textbf{角色} & \textbf{主要职责} \\
\midrule
张可凡（队长） & 全栈开发/项目管理/测试 & Docker 容器化部署、管理后台前端、Postman 接口测试、文档整合 \\
\hline
阮崇轩 & 前端开发 & 前端工程搭建、用户与认证模块（登录/注册/找回密码）、消息通知系统、订单详情与操作、UI 动效与交互优化 \\
\hline
李硕 & 前端开发 & 首页与商品模块（列表/详情/发布）、购物车与订单列表、管理后台页面、图片与数据优化 \\
\hline
刘轩霆 & 后端开发 A & 用户模块（注册/登录/JWT）、商品模块（发布/管理/搜索）、安全框架、订单核心流程、后台管理 API \\
\hline
胡欣悦 & 后端开发 B & 购物车模块、消息与通知模块、数据统计、文件上传、Docker 部署运维 \\
\bottomrule
\end{tabularx}
\end{table}

\newpage
\section{每日工作记录}

\subsection{2026年7月4日（周六）—— 项目启动与环境搭建}

\textbf{全员晨会：} 召开项目启动会，明确技术栈（SpringBoot 3 + Vue 3 + MySQL + Redis）、分工与 Git 工作流。

\begin{table}[H]
\centering
\caption{7月4日工作记录}
\begin{tabularx}{\textwidth}{c X c}
\toprule
\textbf{成员} & \textbf{今日工作内容} & \textbf{完成情况} \\
\midrule
张可凡 & 创建 GitHub 仓库并初始化目录结构；编写 docker-compose.yml 编排 MySQL+Redis+后端+前端容器；配置 Nginx；编写 README.md & $\checkmark$ \\
\hline
阮崇轩 & 使用 Vite 创建 Vue3 + TypeScript 项目；集成 Element Plus、Pinia、Vue Router、Axios；封装 Axios 拦截器；配置路由结构 & $\checkmark$ \\
\hline
刘轩霆 & 协助搭建后端骨架；编写 MyBatis-Plus 配置与分页插件；设计全部 15 张表的建表 SQL 与索引脚本 & $\checkmark$ \\
\hline
胡欣悦 & 搭建 SpringBoot 项目骨架，配置 pom.xml（SpringBoot 3.2.0、MyBatis-Plus 3.5.5、SpringSecurity、JWT 0.12.3）；完成 Maven 构建验证；配置 application.yml 多环境文件 & $\checkmark$ \\
\hline
李硕 & 协助前端项目搭建；完成全部数据库表设计；开发 MainLayout 与 AdminLayout 布局组件 & $\checkmark$ \\
\bottomrule
\end{tabularx}
\end{table}

\textbf{问题：} 无。

\subsection{2026年7月5日（周日）—— 基础框架与核心后端}

\begin{table}[H]
\centering
\caption{7月5日工作记录}
\begin{tabularx}{\textwidth}{c X c}
\toprule
\textbf{成员} & \textbf{今日工作内容} & \textbf{完成情况} \\
\midrule
张可凡 & 完成 Docker 全栈环境联调验证；编写 Redis 配置与缓存工具类；编写 MySQL 与 Redis 配置文件 & $\checkmark$ \\
\hline
阮崇轩 & 开发登录/注册页面 UI（表单校验+验证码输入）；开发 MainLayout 布局；完成 Axios 与后端联调测试 & $\checkmark$ \\
\hline
刘轩霆 & 设计消息表、通知表、公告表 SQL；完成地址管理 5 个接口；完成购物车 4 个接口 & $\checkmark$ \\
\hline
胡欣悦 & 完成统一返回封装、全局异常处理、CORS 配置；完成 SpringSecurity+JWT 认证过滤器；实现注册、登录、验证码三个接口 & $\checkmark$ \\
\hline
李硕 & 开发首页（搜索栏+分类导航卡片+推荐商品网格+公告栏）；开发 AdminLayout 布局 & $\checkmark$ \\
\bottomrule
\end{tabularx}
\end{table}

\textbf{问题：} Docker 容器启动顺序 —— 通过 depends\_on + 健康检查解决。

\subsection{2026年7月6日（周一）—— 用户模块 + 商品模块}

\begin{table}[H]
\centering
\caption{7月6日工作记录}
\begin{tabularx}{\textwidth}{c X c}
\toprule
\textbf{成员} & \textbf{今日工作内容} & \textbf{完成情况} \\
\midrule
张可凡 & 编写需求分析文档；编写系统架构设计文档；编写 migration-indexes.sql 索引优化脚本 & $\checkmark$ \\
\hline
阮崇轩 & 开发个人中心页面（信息编辑+头像上传预览裁剪）；开发地址管理页面；联调用户模块接口 & $\checkmark$ \\
\hline
刘轩霆 & 实现创建订单（事务+库存扣减）；实现订单列表、订单详情、取消订单、订单日志接口 & $\checkmark$ \\
\hline
胡欣悦 & 实现个人信息、头像上传、Token 刷新/退出接口；实现地址管理 5 个接口；实现商品发布、图片上传接口 & $\checkmark$ \\
\hline
李硕 & 开发商品列表页（分类侧栏、卡片网格、分页、排序）；封装 ProductCard.vue；开发商品详情页 & $\checkmark$ \\
\bottomrule
\end{tabularx}
\end{table}

\textbf{问题：} 头像上传需配置静态资源映射 —— 在 WebMvcConfig 中解决。

\subsection{2026年7月7日（周二）—— 订单模块 + 商品管理}

\begin{table}[H]
\centering
\caption{7月7日工作记录}
\begin{tabularx}{\textwidth}{c X c}
\toprule
\textbf{成员} & \textbf{今日工作内容} & \textbf{完成情况} \\
\midrule
张可凡 & 完成数据库设计说明书；开始编写接口设计说明书；开发管理后台用户管理页 & $\checkmark$ \\
\hline
阮崇轩 & 开发首页推荐轮播；开发我的收藏页；开发消息通知页；联调用户模块全部接口 & $\checkmark$ \\
\hline
刘轩霆 & 实现支付（模拟）、确认发货、确认收货、评价接口 & $\checkmark$ \\
\hline
胡欣悦 & 实现商品列表（分页+多条件搜索+排序）、商品详情、编辑/删除/改状态；实现分类列表；实现收藏/取消收藏/收藏列表 & $\checkmark$ \\
\hline
李硕 & 开发商品发布页（多图片上传组件+分类选择器+富文本+价格/成色表单）；联调商品接口 & $\checkmark$ \\
\bottomrule
\end{tabularx}
\end{table}

\textbf{问题：} 多条件组合查询 —— 使用 MyBatis-Plus LambdaQueryWrapper 解决。

\subsection{2026年7月8日（周三）—— 消息模块 + 后台管理}

\begin{table}[H]
\centering
\caption{7月8日工作记录}
\begin{tabularx}{\textwidth}{c X c}
\toprule
\textbf{成员} & \textbf{今日工作内容} & \textbf{完成情况} \\
\midrule
张可凡 & 完成接口设计说明书；开发管理后台商品审核页；开发管理后台订单管理页 & $\checkmark$ \\
\hline
阮崇轩 & 开发站内信页面（会话列表+聊天窗口+发送消息）；联调商品模块接口；编写前端单元测试 & $\checkmark$ \\
\hline
刘轩霆 & 实现站内信 4 个接口；实现系统通知 2 个接口；实现公告 CRUD 3 个接口；实现公开公告列表接口 & $\checkmark$ \\
\hline
胡欣悦 & 实现首页推荐（Redis 缓存热点商品）；编写 JUnit 单元测试；实现分类数据初始化 & $\checkmark$ \\
\hline
李硕 & 开发购物车页；开发订单列表页；开发订单详情页 & $\checkmark$ \\
\bottomrule
\end{tabularx}
\end{table}

\textbf{问题：} 订单库存扣减原子性 —— 通过 @Transactional 事务注解解决。

\subsection{2026年7月9日（周四）—— 管理后台 + 联调}

\begin{table}[H]
\centering
\caption{7月9日工作记录}
\begin{tabularx}{\textwidth}{c X c}
\toprule
\textbf{成员} & \textbf{今日工作内容} & \textbf{完成情况} \\
\midrule
张可凡 & 开发管理后台公告管理页；开发数据统计页（ECharts 图表）；Postman 全量接口测试 & $\checkmark$ \\
\hline
阮崇轩 & 完成站内信页面联调；完成订单页面联调；完成购物车联调；修复前端 Bug & $\checkmark$ \\
\hline
刘轩霆 & 实现后台管理 API（订单管理、数据统计）；编写全部后端单元测试 & $\checkmark$ \\
\hline
胡欣悦 & 实现后台管理 API（用户列表、禁用/启用/改角色、商品审核）；配合前端全量联调 & $\checkmark$ \\
\hline
李硕 & 开发管理后台各页面（用户管理、商品审核、订单管理、公告管理）；集成 ECharts 图表 & $\checkmark$ \\
\bottomrule
\end{tabularx}
\end{table}

\textbf{问题：} 字段命名前后端不一致（下划线 vs 驼峰）—— 统一配置 map-underscore-to-camel-case: true 解决。

\subsection{2026年7月10日（周五）—— 收尾与文档}

\begin{table}[H]
\centering
\caption{7月10日工作记录}
\begin{tabularx}{\textwidth}{c X c}
\toprule
\textbf{成员} & \textbf{今日工作内容} & \textbf{完成情况} \\
\midrule
张可凡 & 编写测试报告文档；编写用户使用说明书；Postman 回归测试；汇总全组文档；编写 README 最终版 & $\checkmark$ \\
\hline
阮崇轩 & 运行全部前端单元测试（8 个），全部通过；优化首页加载性能；修复 UI 兼容性问题 & $\checkmark$ \\
\hline
刘轩霆 & 完成后台管理 API 剩余接口；验证订单状态流转正确性；验证消息收发完整性；协助文档编写 & $\checkmark$ \\
\hline
胡欣悦 & 运行全部后端单元测试（14 个），全部通过；修复边界问题；优化接口响应速度与异常处理 & $\checkmark$ \\
\hline
李硕 & 完成全部管理后台页面开发与联调；验证商城端全流程；运行全部测试并通过 & $\checkmark$ \\
\bottomrule
\end{tabularx}
\end{table}

\textbf{问题：} 无。

\newpage
\section{项目总结}

\subsection{总体完成情况}

\begin{table}[H]
\centering
\caption{项目完成情况一览}
\begin{tabularx}{\textwidth}{l c X}
\toprule
\textbf{模块} & \textbf{完成度} & \textbf{说明} \\
\midrule
环境搭建与部署 & 100\% & Docker Compose 四容器编排，一键启动 \\
用户模块 & 100\% & 注册/登录/验证码/个人信息/头像/地址/Token 共 12 个接口 \\
商品模块 & 100\% & 发布/图片/列表/详情/搜索/排序/分类/收藏共 10 个接口 \\
购物车模块 & 100\% & 增删改查 + 选中结算共 4 个接口 \\
订单模块 & 100\% & 创建/列表/详情/取消/日志/支付/发货/收货/评价共 9 个接口 \\
消息模块 & 100\% & 站内信 4 接口 + 系统通知 2 接口 + 公告 4 接口 \\
后台管理 & 100\% & 用户管理/商品审核/订单管理/数据统计/公告管理共 8 个接口 \\
前端页面 & 100\% & 15+ 个页面，覆盖全部功能 \\
文档 & 100\% & 需求分析/架构设计/数据库设计/接口设计/测试报告/用户手册 \\
测试 & 100\% & 后端 14 个单元测试 + 前端 8 个单元测试 + Postman 全量测试 \\
\bottomrule
\end{tabularx}
\end{table}

\subsection{Git 提交统计}

\begin{table}[H]
\centering
\caption{Git 提交统计}
\begin{tabularx}{\textwidth}{l c c}
\toprule
\textbf{成员} & \textbf{提交次数} & \textbf{新增代码行数} \\
\midrule
张可凡 & 18 & $\sim$2\,500 \\
阮崇轩 & 22 & $\sim$4\,500 \\
刘轩霆 & 20 & $\sim$4\,000 \\
胡欣悦 & 16 & $\sim$3\,500 \\
李硕 & 15 & $\sim$3\,200 \\
\midrule
\textbf{合计} & \textbf{91} & \textbf{$\sim$17\,700} \\
\bottomrule
\end{tabularx}
\end{table}

\subsection{项目交付物清单}

\begin{table}[H]
\centering
\caption{项目交付物}
\begin{tabularx}{\textwidth}{c X}
\toprule
\textbf{类别} & \textbf{交付物} \\
\midrule
后端 & SpringBoot 项目（23+ RESTful API）、SpringSecurity+JWT 安全框架、MyBatis-Plus ORM、Redis 缓存、15 张数据库表 \\
\hline
前端 & Vue3 + TypeScript 项目、15+ 页面、Element Plus 组件库、Pinia 状态管理、ECharts 数据图表 \\
\hline
部署 & Docker Compose 编排（4 容器）、Nginx 反向代理、MySQL/Redis 配置 \\
\hline
文档 & 软件需求规格说明书、系统架构设计说明书、数据库设计说明书、接口设计说明书、测试报告、用户使用说明书 \\
\bottomrule
\end{tabularx}
\end{table}

\end{document}
